\section{Use Cases: Medikamentenverwaltung}

Die folgenden Tabellen beschreiben die funktionalen Anforderungen an die Medikamentenverwaltung. Dabei wird insbesondere die Beziehung zwischen Diagnosen, Medikationen und Arzneimitteln berücksichtigt, um die Integrität des Datenmodells sicherzustellen.

% Use Case 1
\begin{table}[H]
\centering
\renewcommand{\arraystretch}{1.4}
\begin{tabular}{|p{4cm}|p{11cm}|}
\hline
\textbf{Titel} & \textbf{Anlage einer Medikation auf Basis einer Diagnose} \\
\hline
\textbf{Kontext} &
Die Erstellung einer Medikation wird durch einen Arzt ausgelöst, nachdem eine Diagnose für einen Patienten gestellt wurde. Voraussetzung ist, dass sowohl die Diagnose im System existiert als auch ein geeignetes Medikament in der Datenbank verfügbar ist. \\
\hline
\textbf{Beschreibung} &
Zur Behandlung einer diagnostizierten Erkrankung wird eine Medikation angelegt, indem ein konkretes Medikament ausgewählt und mit einer definierten Dosierung sowie Einnahmefrequenz versehen wird. Zusätzlich wird der Beginn der Behandlung dokumentiert, sodass eine zeitliche Nachvollziehbarkeit gewährleistet ist. \\
\hline
\textbf{Besonderheiten} &
Die Medikation muss eindeutig einer bestehenden Diagnose zugeordnet sein. Zudem darf nur auf Medikamente zurückgegriffen werden, die im System als verfügbar geführt werden. \\
\hline
\end{tabular}
\caption{Use Case: Medikation anlegen}
\end{table}

% Use Case 2
\begin{table}[H]
\centering
\renewcommand{\arraystretch}{1.4}
\begin{tabular}{|p{4cm}|p{11cm}|}
\hline
\textbf{Titel} & \textbf{Anpassung einer bestehenden Medikation} \\
\hline
\textbf{Kontext} &
Eine Anpassung erfolgt durch einen Arzt während eines laufenden Behandlungsprozesses, beispielsweise aufgrund veränderter Symptome oder neuer medizinischer Erkenntnisse. Voraussetzung ist eine bereits bestehende aktive Medikation. \\
\hline
\textbf{Beschreibung} &
Eine vorhandene Medikation wird hinsichtlich Dosierung, Medikament oder Einnahmefrequenz angepasst. Ziel ist es, die Behandlung optimal an den aktuellen Zustand des Patienten anzupassen. \\
\hline
\textbf{Besonderheiten} &
Änderungen müssen nachvollziehbar bleiben, sodass der Behandlungsverlauf weiterhin rekonstruierbar ist. Inkonsistente Zustände, wie parallele widersprüchliche Medikationen, sind zu vermeiden. \\
\hline
\end{tabular}
\caption{Use Case: Medikation ändern}
\end{table}

% Use Case 3
\begin{table}[H]
\centering
\renewcommand{\arraystretch}{1.4}
\begin{tabular}{|p{4cm}|p{11cm}|}
\hline
\textbf{Titel} & \textbf{Beendigung einer Medikation} \\
\hline
\textbf{Kontext} &
Die Beendigung wird durch einen Arzt ausgelöst, sobald eine Behandlung abgeschlossen ist oder ein Medikament nicht weiter verabreicht werden soll. Voraussetzung ist eine aktive Medikation ohne gesetztes Enddatum. \\
\hline
\textbf{Beschreibung} &
Die Medikation wird abgeschlossen, indem ein Endzeitpunkt dokumentiert wird. Dadurch wird der Zeitraum der Behandlung vollständig erfasst. \\
\hline
\textbf{Besonderheiten} &
Beendete Medikationen dürfen nicht gelöscht werden, da sie für die medizinische Historie und Nachvollziehbarkeit weiterhin relevant sind. \\
\hline
\end{tabular}
\caption{Use Case: Medikation beenden}
\end{table}

% Use Case 4
\begin{table}[H]
\centering
\renewcommand{\arraystretch}{1.4}
\begin{tabular}{|p{4cm}|p{11cm}|}
\hline
\textbf{Titel} & \textbf{Verwaltung von Arzneimitteln} \\
\hline
\textbf{Kontext} &
Dieser Use Case wird durch administratives Personal oder die Krankenhausapotheke ausgelöst, wenn neue Medikamente eingeführt oder bestehende Einträge aktualisiert werden müssen. Voraussetzung sind entsprechende Zugriffsrechte. \\
\hline
\textbf{Beschreibung} &
Arzneimittel werden im System angelegt, bearbeitet oder deaktiviert. Dabei werden relevante Informationen wie Name, Wirkstoff, Typ sowie der aktuelle Lagerbestand gepflegt. \\
\hline
\textbf{Besonderheiten} &
Nur aktive Medikamente dürfen für neue Medikationen verwendet werden. Änderungen am Bestand müssen konsistent gehalten werden, um eine korrekte Verfügbarkeitsprüfung zu ermöglichen. \\
\hline
\end{tabular}
\caption{Use Case: Arzneimittel verwalten}
\end{table}

% Use Case 5
%\begin{table}[H]
%\centering
%\renewcommand{\arraystretch}{1.4}
%\begin{tabular}{|p{4cm}|p{11cm}|}
%\hline
%\textbf{Titel} & \textbf{Prüfung der Verfügbarkeit eines Medikaments} \\
%\hline
%\textbf{Kontext} &
%Vor der Verschreibung eines Medikaments wird durch medizinisches Personal geprüft, ob dieses im Krankenhaus verfügbar ist. Voraussetzung ist, dass das Medikament im System erfasst ist. \\ 
%\hline
%\textbf{Beschreibung} &
%Das System überprüft den aktuellen Lagerbestand eines Medikaments und stellt fest, ob dieses in ausreichender Menge vorhanden ist oder eine Nachbestellung erforderlich wird. \\
%\hline
%\textbf{Besonderheiten} &
%Die Verfügbarkeitsprüfung ist eine zentrale Voraussetzung für die sichere Versorgung und kann Einfluss auf die Auswahl alternativer Medikamente haben. \\
%\hline
%\end{tabular}
%\caption{Use Case: Verfügbarkeit prüfen}
%\end{table}

% Use Case 6
%\begin{table}[H]
%\centering
%\renewcommand{\arraystretch}{1.4}
%\begin{tabular}{|p{4cm}|p{11cm}|}
%\hline
%\textbf{Titel} & \textbf{Einsicht in die Medikationshistorie eines Patienten} \\
%\hline
%\textbf{Kontext} &
%Dieser Use Case wird durch Ärzte oder Pflegepersonal ausgelöst, wenn Informationen über vergangene oder aktuelle Medikationen benötigt werden. Voraussetzung ist, dass entsprechende Daten im System vorhanden sind. \\
%\hline
%\textbf{Beschreibung} &
%Alle bisherigen Medikationen eines Patienten werden aus dem System geladen und chronologisch dargestellt, sodass ein vollständiger Überblick über den Behandlungsverlauf entsteht. \\
%\hline
%\textbf{Besonderheiten} &
%Die Historie dient als Grundlage für medizinische Entscheidungen und zur Vermeidung von Wechselwirkungen zwischen Medikamenten. \\
%\hline
%\end{tabular}
%\caption{Use Case: Medikationshistorie einsehen}
%\end{table}

\vspace{1em}

\subsection*{Resultierende Entitäten}

\begin{itemize}
\item \textbf{Diagnosis (Diagnose):} Repräsentiert die medizinische Feststellung einer Erkrankung und bildet die Grundlage für therapeutische Maßnahmen.
\item \textbf{Medication (Medikation):} Verknüpft eine Diagnose mit einem konkreten Medikament sowie zeitlichen und dosierungsbezogenen Informationen.
\item \textbf{Drugs (Arzneimittel):} Enthält alle verfügbaren Medikamente inklusive Eigenschaften wie Name, Wirkstoff und Lagerbestand.
\item \textbf{Dosis (Dosierung):} Definiert Menge, Einheit und Einnahmefrequenz und konkretisiert damit die Medikation.
\end{itemize}
